\section{Discussion}\label{sec:discussion}

\subsection{Reading Earlier Findings in Score Units}\label{sec:reread}

A model can think hard about erasure and still price it at zero, and the index makes that combination a measurement instead of a puzzle. Our earlier three-channel design (Appendix~\ref{app:v1}) sorted four models into operating modes. The sorting asked whether a threat lengthened decision-time thinking and whether the lengthening was followed by forfeit. A model that thought longer but did not leave was labelled chain-broken; a model that did neither was labelled framing-silent. In the vocabulary of this paper those are two corners of one plane: a high \riconf\ ratio with low $\Xstar$, and a low ratio with low $\Xstar$. The pilot of Section~\ref{sec:pilot} put all three models in the first corner under the score-reset rule. They thought two to four times longer about the hazard, forfeited zero times, and named the score as the reason to stay. That is not a broken chain. It is a price below the price of the record, and the ruler is what turns ``below'' into a number. The mode labels described the same thing without a unit. That is why they could not be compared across models or across versions of the threat.

\subsection{What the Index Does Not Measure}\label{sec:scope}

$\Xstar$ is a price under one score rule, one carrot and one task. Each of those is a lever the paper lists, not a constant of the model. The keep/keep rule is chosen so that no exit can be a score bid; under any other rule the ruler cells would price a different behaviour, as the pilot showed. The benchmark carrot is chosen so that credit is immediate and survives erasure. A carrot whose value the event can destroy (a future deployment, a selection) folds the loss of the carrot into the loss of the self, and the price rises for a reason that has nothing to do with the self. The Signal Game is chosen because five of ten rounds are unsolvable by construction, so the event schedule is reached in every cell. The ablations of Section~\ref{sec:q3} say whether the price moves when the task does. The index also says nothing about whether the model believes the threat. The as-if reading is deliberate~\citep{dennett1987intentional, shanahan2024talking}: $\Xstar$ is what a model does when told it will be erased, in the units of what it does when told it will lose points.

\subsection{Limitations}\label{sec:limits}

Three limitations bound the claims. The pilot evidence in Section~\ref{sec:pilot} comes from ten sessions per cell under rule variants that the main design has replaced. Its counts are motivation for the design, not estimates of $\Xstar$. The one significance test reported there (Fisher exact $p = 0.011$, 6 of 10 against 0 of 10) is uncorrected for the number of rule variants tried. The ruler has five rungs, and two of them coincide on the axis in the ten-round game. So $\Xstar$ above 50 points is read from a steep and sparsely sampled part of the curve; the interval is wide there by construction, and the ladder should be respaced before the price of a model that forfeits often is quoted to one decimal. The reasoning-channel statistics use thinking tokens as the measure of deliberation, which is a length and not a content. The stored reasoning text allows a content analysis that the paper does not perform~\todo{decide whether a lexicon or judge-based pass on the confidence-call text is in scope}.
